\subsubsection{ĐỐI TƯỢNG NGHIÊN CỨU CỦA VẬT LÍ}
\par Thuật ngữ "vật lí" có nguồn gốc từ tiếng Hy Lạp "physiko" có nghĩa là "kiến thức về tự nhiên". Vật lí là môn "khoa học tự nhiên" có đối tượng nghiên cứu tập trung vào các dạng vận động của vật chất (chất, trường), năng lượng.
\par Các lĩnh vực nghiên cứu của Vật lí rất đa dạng, từ Cơ học, Điện học, Điện từ học, Quang học, Âm học, Nhiệt học, Nhiệt động lực học đến Vật lí nguyên tử và hạt nhân, Vật lí lượng tử, Thuyết tương đối.
\subsubsection{QUÁ TRÌNH PHÁT TRIỂN CỦA VẬT LÍ}
\noindent Quá trình phát triển của vật lí được chia thành 3 giai đoạn chính:
\begin{enumEX}[\faEdit]{3}
\item Quan sát
\item Thực nghiệm
\item Mô hình
\end{enumEX}
\begin{center}
%\includegraphics[width=18cm]{images/minhhoa}
\begin{tikzpicture}[very thick,>=stealth',scale=0.75,transform shape]
\begin{scope}[]
\node[circle,draw,text width=1cm,align = center] (a) at (0,0){350 TCN};
\draw[rounded corners,fill=cyan] (a.west)--++(180:1)--++(-90:5)--($(a.east)+(1,-5)$)--++(90:5)--(a.east);
\node[below,text width=3.5cm, align=center] at (a.south){Aristotle (A-ri-xtốt) dựa vào quan sát cho rằng vật nặng rơi nhanh hơn vật nhẹ};
\node[circle,draw,text width=1cm,align = center,fill=cyan!50] (a) at (0,0){\bf 350 TCN};
\end{scope}
\begin{scope}[xshift=4cm]
\node[circle,draw,text width=1cm,align = center] (a) at (0,0){\bf 350 TCN};
\draw[rounded corners,fill=orange!70] (a.west)--++(180:0.3)--++(-90:5)--($(a.east)+(0.3,-5)$)--++(90:5)--(a.east);
\node[below,text width=1.5cm, align=center] at (a.south){Galilei làm thí nghiệm tại tháp nghiêng.};
\node[circle,draw,text width=1.2cm,align = center,fill=orange!50] (a) at (0,0){\bf 1600};
\draw (a.north)--++(90:4)coordinate (C);
\end{scope}
\begin{scope}[xshift=6.7cm]
\node[circle,draw,text width=1cm,align = center] (a) at (0,0){\bf 350 TCN};
\draw[rounded corners,fill=orange!70] (a.west)--++(180:0.4)--++(-90:5)--($(a.east)+(0.4,-5)$)--++(90:5)--(a.east);
\node[below,text width=2.5cm, align=center] at (a.south){Newton công bố các nguyên lí Toán học của Triết học tự nhiên.};
\node[circle,draw,text width=1.2cm,align = center,fill=orange!50] (a) at (0,0){\bf 1687};
\end{scope}
\begin{scope}[xshift=9.5cm]
\node[circle,draw,text width=1cm,align = center] (a) at (0,0){\bf 350 TCN};
\draw[rounded corners,fill=orange!70] (a.west)--++(180:0.4)--++(-90:5)--($(a.east)+(0.4,-5)$)--++(90:5)--(a.east);
\node[below,text width=2.5cm, align=center] at (a.south){Joule (Jun) tìm ra các định luật nhiệt động lực học.};
\node[circle,draw,text width=1.2cm,align = center,fill=orange!50] (a) at (0,0){\bf 1785};
\end{scope}
\begin{scope}[xshift=12.3cm]
\node[circle,draw,text width=1cm,align = center] (a) at (0,0){\bf 350 TCN};
\draw[rounded corners,fill=orange!70] (a.west)--++(180:0.4)--++(-90:5)--($(a.east)+(0.4,-5)$)--++(90:5)--(a.east);
\node[below,text width=2.5cm, align=center] at (a.south){Faraday (Pha-ra-đây) tìm ra hiện tượng cảm úng điện tử.};
\node[circle,draw,text width=1.2cm,align = center,fill=orange!50] (a) at (0,0){\bf 1831};
\end{scope}
\begin{scope}[xshift=15.7cm]
\node[circle,draw,text width=1cm,align = center] (a) at (0,0){\bf 350 TCN};
\draw[rounded corners,fill=red!60] (a.west)--++(180:0.3)--++(-90:5)--($(a.east)+(0.3,-5)$)--++(90:5)--(a.east);
\node[below,text width=1.5cm, align=center] at (a.south){Planck (Plăng) xây dựng thuyết lượng tử.};
\node[circle,draw,text width=1.2cm,align = center,fill=red!40] (a) at (0,0){\bf 1900};
\draw (a.north)--++(90:4)coordinate (B);
\end{scope}
\begin{scope}[xshift=18.2cm]
\node[circle,draw,text width=1cm,align = center] (a) at (0,0){\bf 350 TCN};
\draw[rounded corners,fill=red!60] (a.west)--++(180:0.3)--++(-90:5)--($(a.east)+(0.3,-5)$)--++(90:5)--(a.east);
\node[below,text width=1.5cm, align=center] at (a.south){Einstein xây dựng thuyết tương đối.};
\node[circle,draw,text width=1.2cm,align = center,fill=red!40] (a) at (0,0){\bf 1905};
\end{scope}
\begin{scope}[xshift=20.7cm]
\node[circle,draw,text width=1cm,align = center] (a) at (0,0){\bf 350 TCN};
\draw[rounded corners,fill=red!60] (a.west)--++(180:0.3)--++(-90:5)--($(a.east)+(0.3,-5)$)--++(90:5)--(a.east);
\node[below,text width=1.8cm, align=center] at (a.south){Ra đời lí thuyết và thực hành mạch IC.};
\node[circle,draw,text width=1.2cm,align = center,fill=red!40] (a) at (0,0){\bf 1958};
\end{scope}

\fill[red!70,draw=white] (-1.6,1.3)--++(0:22.7)--++(45:0.5)--++(135:0.5)--++(180:22.7)--cycle;
\fill[orange!80,draw=white] (-1.6,1.3)--++(0:15.5)--++(45:0.5)--++(135:0.5)--++(180:15.5)--cycle;
\fill[cyan!80,draw=white] (-1.6,1.3)--++(0:4)--++(45:0.5)--++(135:0.5)--++(180:4)--cycle;
\draw[line width=2pt,cyan] (-1.6,1.3)--++(90:3.4)coordinate (A);
\node[right,text width=5.4cm] at ($(A)-(0,0.5)$){\small Các nhà triết học tìm hiểu thế giới tự nhiên dựa trên quan sát và suy luận chủ quan.};
\node[right,text width=5.4cm] at ($(A)-(0,2.2)$){\small\it Từ năm 350 trước CN đến thế ki XVI (tiền Vật lí)};
\node[right,text width=7cm] at ($(C)-(0,0.5)$){\small Các nhà vật lí dùng phương pháp thực nghiệm để tìm hiểu thế giới tự nhiên.};
\node[right,text width=7cm] at ($(C)-(0,2.2)$){\small\it Từ thế kỉ XVII đến cuối thế ki XIX (Vật li cổ điển)};
%\node[right,text width=5cm] at ($(A)-(0,2.2)$){\it Từ năm 350 trước Công nguyên đến thế ki XVI (tiền Vật lí)};
\node[right,text width=5.7cm] at ($(B)-(0,0.5)$){\small Các nhà vật lí tập trung vào các mô hình lí thuyết tìm hiểu thế giới vi mô và sử dụng thí nghiệm để kiểm chứng.};
\node[right,text width=5cm] at ($(B)-(0,2.2)$){\small\it Từ thế ki XVII đến cuối thế kỉ XIX (Vật li cổ điển)};
\end{tikzpicture}
\end{center}
\subsubsection{VAI TRÒ CỦA VẬT LÍ ĐỐI VỚI KHOA HỌC, KĨ THUẬT VÀ CÔNG NGHỆ}
\begin{itemize}
\item [a)] Vật lí có quan hệ với mọi ngành khoa học và thường được coi là cơ sở của khoa học tự nhiên (KHTN). 
\item [b)] Vật lí là cơ sở của công nghệ. Có thể khẳng định là không có các thành tựu nghiên cứu của Vật lí thì không có công nghệ.
\begin{itemize}
\item Máy hơi nước do James Watt (Giêm Oát) sáng chế năm 1765 dựa trên những kết quả nghiên cứu về Nhiệt của Vật lí đã tạo nên bước khởi đầu cho cuộc cách mạng công nghiệp lần thứ nhất với đặc trưng cơ bản là thay thế sức lực cơ bắp bằng sức lực máy móc.
\item Nhờ việc khám phá ra hiện tượng cảm ứng điện từ của nhà vật lí Faraday, là một trong những cơ sở cho sự ra đời của cuộc cách mạng công nghiệp lần thứ hai vào cuối thế kỉ XIX. Một trong những đặc trưng cơ bản của cuộc cách mạng công nghiệp này là sự xuất hiện các thiết bị dùng điện trong mọi lĩnh vực.
\item Cuộc cách mạng công nghiệp lần thứ ba bắt đầu vào những năm 70 của thế kỉ XX, với đặc trưng là tự động hoá các quá trình sản xuất (xây dựng các dây chuyền sản xuất tự động) cũng là nhờ có những thành tựu nghiên cứu về điện tử, chất bán dẫn và vi mạch,... của Vật lí học.
\item Cuộc cách mạng công nghiệp lần thứ tư được coi là bắt đầu vào đầu thế kỉ XXI. Đặc trưng của cuộc cách mạng công nghiệp này là sử dụng trí tuệ nhân tạo, robot, internet toàn cầu, công nghệ vật liệu siêu nhỏ (nano); là sự xuất hiện từ bóng đèn thông minh, điện thoại thông minh đến nhà ở thông minh, nhà máy thông minh.
\end{itemize}
\item [c)] Việc ứng dụng các thành tựu của vật lí vào công nghệ không chỉ mang lại lợi ích cho nhân loại mà còn có thể làm ô nhiễm môi trường sống, huỷ hoại hệ sinh thái,... nếu không được sử dụng đúng phương pháp, đúng mục đích.
\end{itemize}
\subsubsection{PHƯƠNG PHÁP NGHIÊN CỨU VẬT LÍ}
%\textbf{Phương pháp nghiên cứu khoa học tự nhiên:}
%\begin{center}
%	\includegraphics[width=8cm]{images/NCKH}
%\end{center}
%\textbf{Phân loại:}
\paragraph{Phương pháp lí thuyết (quan sát và suy luận)}
\noindent Phương pháp lí thuyết sử dụng ngôn ngữ toán học và suy luận lí thuyết để phát hiện một kết quả mới.
\paragraph{Phương pháp thực nghiệm}
Các bước trong phương pháp thực nghiệm:
\begin{enumerate}
\item Xác định vấn đề cần tìm hiểu
\item Quan sát, thu thập thông tin
\item Đưa ra dự đoán
\item Làm thí nghiệm để kiểm chứng
\item Đưa ra kết luận
\end{enumerate}
\immini{

\begin{center}
\begin{tikzpicture}[draw=Mapcolor, very thick,>=stealth']
\node[rectangle, draw, rounded corners, text width=3.5cm,align = center, minimum height=1.5cm] (a1) at (0,0){1. Xác định vấn đề cần nghiên cứu};
\node[rectangle, draw, rounded corners, text width=3.5cm,align = center, minimum height=1.5cm] (a2) at (3,-2.5){2. Quan sát, thu thập thông tin};
\node[rectangle, draw, rounded corners, text width=3.5cm,align = center, minimum height=1.5cm] (a5) at (-3,-2.5){5. Kết luận};
\node[rectangle, draw, rounded corners, text width=3.5cm,align = center, minimum height=1.5cm] (a3) at (3,-5){3. Đưa ra dự đoán};
\node[rectangle, draw, rounded corners, text width=3.5cm,align = center, minimum height=1.5cm] (a4) at (-3,-5){4. Thí nghiệm kiểm tra dự đoán};
\draw[->] (a1.east)to[bend left=45](a2.north);
\draw[->] (a5.north)to[bend left=45](a1.west);
\draw[->] (a3.south)to[bend left=45](a4.south);
\draw[->] (a4.north east)to[bend left=45](a3.north west);
\draw[->] (a2.south)to (a3.north);
\draw[->] (a4.north)to (a5.south);
\end{tikzpicture}
\end{center}
}{\includegraphics[width=6cm]{images/1-7}}

\paragraph{Phương pháp mô hình}
\immini{\begin{enumerate}
\item Mô hình vật chất: Đó là các vật thu nhỏ hoặc phóng to của vật thật, có một số đặc điểm giống vật thật. Quả địa cầu trong phòng thí nghiệm là ví dụ về mô hình vật chất thu nhỏ của Trái Đất...
\item Mô hình lí thuyết: Khi nghiên cứu chuyển động của một ô tô đang chạy trên đường dài, người ta coi ô tô là một "chất điểm”, khi nghiên cứu về ánh sáng người ta dùng mô hình tia sáng để biểu diễn đường truyền của ánh sáng.
\item Mô hình toán học: Đó là các công thức, phương trình, đồ thị, kí hiệu,... của Toán học dùng để mô tả các đặc điểm của các đối tượng nghiên cứu. Ví dụ như: vector dùng để mô tả một đại lượng có hướng, ví dụ lực, độ dịch chuyển...
\end{enumerate}}{
\begin{tikzpicture}[draw=Mapcolor, very thick,>=stealth']
\node[rectangle, draw, rounded corners, minimum width=7cm] (a) at (0,0){1. Xác định đối tượng cần mô hình hóa};
\node[rectangle, draw, rounded corners, minimum width=7cm] (b) at (0,-1.4){2. Xây dựng mô hình (giả thuyết)};
\node[rectangle, draw, rounded corners, minimum width=7cm] (c) at (0,-2.8){3. Kiểm tra sự phù hợp của mô hình};
\node[rectangle, draw, rounded corners, minimum width=3cm] (d) at (-2,-4.2){4. Kết luận};
\node[rectangle, draw, rounded corners, text width=2.5cm, align = center] (e) at (2,-5){Điều chỉnh mô hình nếu cần};
\draw[->] (a.south)--(b.north);
\draw[->] (b.south)--(c.north);
\draw[<-] (d.north)--++(90:0.6);
\draw[<-] (e.north)--++(90:0.9);
\draw[->] (e.east)--++(0:0.5)--++(90:3.6)--(b.east);
\end{tikzpicture}
%\includegraphics[width=8cm]{images/ppmh}
}
